% Partie inférence (Quentin) : rejet par le discriminateur (DRS, OBRS), recalibration de D.
% Sources : enquete/results/juliet/drs_recalib_test.csv (tableau, tests appariés : figscripts/inference.py),
% juliet2_test.csv (G brut), enquete/inference.md (écart-type des logits, AUC G brut / G_EMA, coût).
\section{Inference: filtering images with the discriminator}\label{sec:inference}

\textbf{Idea.} Even a good generator produces some bad images, and we have a network trained precisely to spot them: the
discriminator. Discriminator rejection sampling (DRS)~\cite{azadi2019} uses it as a filter: draw many images, and keep
each one with a probability that grows with $D$'s score. If $D$ were perfect, its score would tell exactly how much more
likely an image is under the real data than under $G$, and the filter would correct $G$'s mistakes. OBRS~\cite{verine2024}
is the optimal version of this filter for a fixed budget: it keeps every image above a threshold and fewer below, with
the threshold set to keep one image out of three. On ImageNet, DRS improved the FID from 18.2 to 17.1~\cite{azadi2019}.

\begin{table}[t]
\centering\small
\setlength{\tabcolsep}{4.5pt}
\caption{\textbf{Filtering with $D$.} 10\,000 kept images, 3 seeds $\times$ 2 draws.}
\label{tab:inference}
\begin{tabular}{@{}lcccc@{}}
\toprule
Filter & AUC of $D$ & Kept & \finc{} & \fcnn{} \\
\midrule
\multicolumn{5}{@{}l}{\emph{Baseline}} \\
none                 & --   & 100\% & $38.7\pm0.9$ & 183 \\
DRS                  & 0.75 & 30\%  & $34.3\pm0.8$ & 154 \\
OBRS                 & 0.75 & 33\%  & $\mathbf{33.8\pm0.9}$ & 145 \\
OBRS, recalibrated   & 0.83 & 33\%  & $36.1\pm0.9$ & 83 \\
\midrule
\multicolumn{5}{@{}l}{\emph{Final model (RpGAN + R1 + EMA)}} \\
none                 & --   & 100\% & $22.0\pm1.4$ & 9.9 \\
DRS                  & 0.57 & 42\%  & $22.1\pm1.1$ & 14.4 \\
DRS, recalibrated    & 0.77 & 32\%  & $20.2\pm1.2$ & 10.0 \\
OBRS, recalibrated   & 0.77 & 33\%  & $\mathbf{19.7\pm0.8}$ & 10.2 \\
\bottomrule
\end{tabular}
\tabnote{AUC: probability that $D$ scores a real image above a generated one (0.5 = coin flip, 1 = perfect).}
\end{table}

\textbf{Result.} On the baseline, the filter works: DRS improves \finc{} from 38.7 to 34.3, OBRS slightly more
(Table~\ref{tab:inference}). On the final model, it does nothing to \finc{} (22.0 to 22.1) and even worsens \fcnn{} on
every seed: this is the degradation we had reported in our presentation.

\textbf{Why it fails: $D$ can no longer tell real from fake.} To measure how well $D$ still separates the two, we use its
AUC (area under the ROC curve): take one real and one generated image at random; the AUC is the probability that $D$ gives
the real one the higher score. A perfect judge scores 1, a coin flip 0.5. The baseline's $D$ reaches 0.75; the final
model's $D$ only 0.57, barely better than chance. This is a side effect of the training that made $G$ good: $G$ and $D$ end near a balance where $D$ is confused. A
filter based on a confused judge keeps images almost at random (42\% instead of the intended 30\%), and the little
preference it has shifts the balance between digits, which only \fcnn{} sees. Part of the confusion comes from $D$ never
having seen the averaged generator, but DRS also fails on the raw generator, so this is not the whole story.

\textbf{Fix: recalibrate $D$.} As Azadi et al.\ did for their model, we keep $D$'s internal features and retrain only its
last layer, a logistic regression fitted on 10\,000 real validation images against 10\,000 generated ones. The AUC rises
to 0.77, and both filters now help on every seed: DRS reaches 20.2 and OBRS 19.7 ($p=0.02$). On the baseline, by
contrast, recalibration rebalances the digits but loses to plain DRS in \finc{}: like the rebalancing of
Sec.~\ref{sec:fid}, it replaces easy, well-drawn ``1''s with harder digits that the baseline draws badly.

\textbf{Cost.} Keeping 30--45\% of images means drawing 2.2 to 3.3 images per kept one. Generating 10\,000 images on one
A100 GPU takes 0.02\,s without filter, 0.08--0.09\,s with DRS and 0.16\,s with OBRS, plus fitting the new layer once.
We found the recalibration after submitting, so the submitted model uses no filter.

\takeaway{Filtering with $D$ only helps if $D$ can still tell real from generated images. A well-trained GAN loses this
ability, but retraining $D$'s last layer restores it and gains about 2 points.}
